\chapter{Introduction}


In image-guided interventions, a clinician must know the exact position of their instruments. Tracking systems solve this by measuring the instrument's position and orientation (pose) continuously. The two main technologies used are optical tracking and electromagnetic tracking (EMT) \cite{sorriento2020}. Optical tracking is accurate and robust, but requires line-of-sight between the markers and a camera. However, electromagnetic tracking does not require any line-of-sight. It instead uses magnetic fields to locate the sensor. This comes at a cost. Metallic and ferromagnetic objects in the workspace distort the field and degrade accuracy \cite{sorriento2020}. 


The Anser EMT system \cite{jaeger2017}, developed at University College Cork (UCC) is one such system. It uses the magnetic field generated by eight coils to estimate the pose of a small sensor near the field. Recovering a pose from eight coil measurements is a nonlinear inverse problem with no closed form. For that reason, numerical methods are used to solve it. The Levenberg-Marquardt (LM) algorithm is the tool used by the Anser system. It is precise but needs a good initialisation to work. When the system is already running, the previous pose can be used as an initialisation for the next but, when the system is first started or loses signal, providing a good first initialisation is very important \cite{jaeger2018}. In practice, this is addressed by placing the sensor at a known pose during startup. However, this solution constrains how the system can be deployed. This thesis asks the question: 'how do we construct a good first initialisation?' 

This thesis constructs a hybrid solver. A neural network (NN) which is imprecise but reliable, provides the initialisation, and Levenberg-Marquardt, which is precise but dependent on a good start, refines it.

Qin et al. \cite{qin2022} previously used a similar hybrid approach for permanent magnet localisation. Three things distinguish this work from theirs. Firstly, the forward map is made up of the integral over many filaments of an eight coil emitter system rather than a dipole approximation. Secondly, the basin of attraction is characterised directly, seperately in position and orientation, rather than assumed from a cited threshold. Thirdly, the network is trained against a loss which weights positional and angular error to equal range, rather than summing them unweighted, a choice which Chapter \ref{chap:results} shows accounts for the largest single improvement in network accuracy. 



This thesis, and its corresponding codebase, provides:
\begin{itemize}
  \item A Python reimplementation of the Anser forward model and solver. 
  \item A characterisation of LM's convergence behaviour as a function of initialisation quality, in both position and orientation.
  \item A feed-forward NN trained to predict pose directly. 
  \item A hybrid solver which uses the network output to provide initialisation for Levenberg-Marquardt.
  \item A three-way empirical comparison establishing that this hybrid solver is both more reliable and faster than LM.
\end{itemize}
\textbf{Chapter \ref{chap:background}} derives the forward map and introduces two solver classes: iterative solvers, and NN. 
\\
\textbf{Chapter \ref{chap:methods}} collects the methods required to build and evaluate the solvers. 
\\
\textbf{Chapter \ref{chap:results}} shows that using NNs to initialise a LM solver both improves reliability and reduces computation time, compared to a cold-start Levenberg-Marquardt solver, without sacrificing precision.
